\section{研究现状、遗留课题与未来方向}\label{sec:future}

\subsection{当前研究现状}

把两条脉络的最新进展放在一起看，有四个趋势值得注意。

\textbf{其一，两条脉络正在方法上合流}。最直接的证据是“慢动量加快反转”这类频段融合工作\cite{wood2021}。更广泛的迹象是机器学习方法的介入：Gu 等(2020) 的资产定价实证显示，树模型与神经网络的收益预测能力显著优于线性模型，且动量类变量在其中扮演重要角色\cite{gu2020}。当模型不再预设“动量或反转”，而是从数据中学习响应函数的形状（甚至允许它随状态变化）时，二者作为\emph{两个候选假设}的地位被削弱，变成\emph{同一个响应函数在不同状态下的取值}。

\textbf{其二，反转一侧正在从“选对”转向“构造组合”}。早期配对交易的核心问题是选对（距离法、协整法）；Cuturi 与 d'Aspremont(2013) 的方差阈值方法把问题改写为：直接求解使均值回归程度最大化、同时约束方差的线性组合\cite{cuturi2013}。这一转向的实质是承认“配对”只是构造均值回归载体的一种启发式，更一般的问题是优化式地合成这样的载体。近年 OU 框架的工程化复现（如 arXiv:2412.12458）与深度学习统计套利（arXiv:2106.04028）延续了这一路线\cite{suchato2024,guijarro2021}；强化学习方案的引入（如 arXiv:2403.12180）则试图把仓位管理与信号生成放进同一学习目标\cite{ning2024}。

\textbf{其三，趋势一侧的重心转向容量与拥挤}。信号构造问题在 2010 年代基本定型（形成期/持有期、波动率缩放、信号平滑），此后文献的注意力明显移向容量与拥挤：Stein(2009) 的理论、Lou--Polk(2022) 的 comomentum 度量、Baltas(2019) 的拥挤综述\cite{stein2009,loupolk2022,baltas2019}。这反映了一个现实：趋势策略的信号收益难以继续提升，制约它的是“钱太多”而非“信号太弱”。

\textbf{其四，两策略都进入了“事后审计”时代}。紧缩夏普比率、回测过拟合概率、多重检验校正已从学术规范变成行业惯例\cite{bailey2014,bailey2016,hlz2016,harveysliu2015}。这对两个策略是同等的要求。

\subsection{遗留课题}

以下课题在本文的证据梳理中反复出现，且都未解决。

\textbf{Q1：U 形自相关是否跨市场普适？} 本文的对照实证给出了一个值得警惕的反例：在中国期货的\emph{时序}口径上，趋势策略在所有 $L$ 上都优于反转策略，“短尺度反转”未出现；U 形结构只出现在\emph{截面}口径上（短 $L$ 趋势占优、中长 $L$ 反转占优）。这说明“尺度决定归属”的普适性依赖具体口径与市场。其成因（投资者结构、涨跌停制度、T+1 交易）尚缺系统研究。

\textbf{Q2：趋势收益中“预测性”与“风险缩放”各占多少？} Kim 等(2016) 指出时间序列动量的收益很大部分来自波动率缩放本身\cite{kim2016}；Huang 等(2020) 发现其表现与简单历史均值规则几乎相同\cite{huang2020}。把这两个成分干净地分解，是判定“趋势溢价是否真实”的关键，但尚无公认的分解方案。

\textbf{Q3：拥挤度能否事前度量？} comomentum 属于滞后代理（从已实现相关性推断），无法事前预警；领先代理（资金流、持仓集中度）数据可得性差。一个可事前计算、且与未来收益衰减确有预测关系的拥挤指标，至今没有。

\textbf{Q4：反转策略的容量上限能否给出定量估计？} 与趋势不同，反转的约束来自自身成本，理论上更可测算（带宽与成本的三次方根关系提供了一个起点\cite{martin2011}）。但从理论带宽到实际可管理规模之间缺少完整的映射。

\textbf{Q5：两策略相关性的状态依赖如何建模？} 二者在某些状态下高度负相关（信号相反），在资金冲击下正相关（同为去杠杆受害者）。这一“相关性变号”现象对组合构建是致命的，但目前多被当作经验事实处理，缺少结构模型。

\textbf{Q6：非线性冲击成本如何改变比较结论？} 本文的成本模型是固定基点。反转策略的交易密集发生在流动性差的时点，其真实冲击成本高于平均值；若采用参与率约束的平方根模型，反转的净收益会被进一步压低。这一敏感性尚未被系统量化。

\textbf{Q7：做空约束下反转策略的可实现性？} 反转的多空结构依赖空头端，而多数新兴市场（含 A 股）做空困难。纯多头版本的反转策略（只买输家）在收益与风险上如何，缺少系统研究。

\textbf{Q8：机器学习是否改变了两策略的本质区分？} 若模型可以自由学习响应函数形状，“趋势”与“反转”退化为两个特例。那么二者的区分是否还具有结构意义，取决于响应函数在实际数据中是否呈现稳定的分段结构——这是一个可检验但尚未被检验的问题。

\textbf{Q9：机制之争能否被判决？} 目前的风险解释与行为解释在价格数据上难以区分\cite{cochrane2011}。一个可能的出路是寻找\emph{只在其中一种机制下成立}的辅助预测（如资金流数据、投资者结构变化），但目前尚无成功案例。

\textbf{Q10：极端资金冲击下两策略同时失效，如何对冲？} 2007 年 8 月的教训是：冲击来自资金端，分散化失效\cite{khandanilo2011,brunnermeier2009}。一个能在资金冲击下真正分散的机制（如显式的流动性保险、跨市场分散、或有减仓规则）尚未被提出。

\subsection{未来方向}

\textbf{方向一：频段分解的统一框架}。把收益序列按时间尺度分解（小波、经验模态分解或频域滤波），分别在每个频段上估计自相关符号与幅度，从而把两策略从“二选一”变成“多频段风险预算”。这一方向在文献中已有雏形\cite{wood2021}，但缺少系统化的理论与实证。它也是本文统一框架最自然的延伸。

\textbf{方向二：状态感知的动态融合}。用 regime 切换模型（Hamilton 1989；Ang--Timmermann 2012）估计“趋势占优”与“反转占优”的状态概率，据此动态分配两策略的权重\cite{hamilton1989,angtimmermann2012}。关键在于状态变量必须\emph{可事前观测}，否则退化为事后标定。

\textbf{方向三：信号—成本—资金的联合优化}。Gârleanu 与 Pedersen(2013) 给出了含成本的最优动态组合解\cite{garleanu2013}，但未纳入资金约束与保证金机制。把三者纳入同一优化问题，对反转策略尤其重要——它的成败往往由资金端决定。

\textbf{方向四：中国市场的制度性比较研究}。中国市场的制度特征（T+1、涨跌停、融券受限、投资者结构以散户为主）为检验“两策略差异的制度依赖性”提供了独特样本。本文的对照实证只是初步观察，其成因（尤其是时序口径下短反转的缺失）值得专门研究。

\textbf{方向五：事前拥挤度量}。从持仓集中度、资金流、期权隐含信息等可得数据构造可事前计算的拥挤代理，并检验其与趋势策略未来收益衰减的关系。这是把第~\ref{sec:duality}~节“容量约束”对偶从定性判断推进到定量工具的关键。

\subsection{全文结论}

本文从一句直觉出发——“趋势与反转只是同一信号的两种符号选择”——做了三件事：把两条脉络编年并置，给出统一与对偶的数学框架，并用统一口径的对照实证检验核心预测。

主要的结论可以压缩成四点。

\textbf{第一，统一是真的，但只统一在信号层}。两策略共享同一信号（过去收益的符号或对均值的偏离），共享同一数学容器（带符号回归强度的 OU 族），共享同一统计量（方差比/自相关函数）。凡涉及“用什么信息、赌什么量”的问题，二者是同一个东西。

\textbf{第二，差异是真的，而且集中在执行层}。平稳性假设、收益偏度、时间尺度、风险来源、容量与换手、失效模式、入场门槛——七组对偶中没有一组是修辞性的，它们都可以追溯到具体的工程后果：反转要拓宽带宽、要做因子剥离、要留资金冗余；趋势要平滑信号、要做状态过滤、要监控拥挤度。

\textbf{第三，成本是两策略最不对称的一把尺子}。本文的对照实证给出了一个干净的示范：在期货时序口径、$L=5$ 的设定下，趋势与反转的换手几乎完全相同（52.41 对 52.38），成本拖累完全相等；但由于二者毛收益互为相反数，成本这一“常数扣项”让表现好的一侧被削去一截、表现差的一侧被推向更差——净夏普分别为 $+0.20$ 与 $-0.60$，而非预期的 $+0.4$ 与 $-0.4$。这不是数据噪声，而是成本对符号的必然作用：\emph{成本不区分方向，它只惩罚交易}。谁更频繁地交易，谁就更多地被惩罚。

\textbf{第四，机制叙事应当被降权}。关于“风险还是行为”的争论无法从价格数据中裁决，它会影响我们对策略稳健性的先验，但不应当成为选择策略的依据。可依据的是可观测的结构事实：自相关在哪个尺度上显著、成本与容量的实测约束、资金端的状态。

回到最初的问题——两个策略的本质异同是什么？本文的答案是：\textbf{它们赌的是同一件事（收益的符号），只是押在了符号的两端；它们的分野不在信号里，而在信号之后的每一步——验证、缩放、对冲、交易、融资}。把这句话说清楚的实际价值是：当有人问“该用趋势还是反转”时，正确的问题不是“你相信市场会走远还是会回头”，而是“你要在哪个尺度上、以什么成本、承担什么样的资金约束，去押注哪个符号”。
